% =============================================================================
\subsection{Real-world evaluation: localization (parking lot) and tracking (ARKit walkaround)}
\label{sec:eval-realworld}
% =============================================================================

\begin{figure}[t]
  \centering
  \includegraphics[width=\linewidth]{figures/combined_localization_tracking.png}
  \caption{Real-world evaluation of Tier~1 + Tier~2 on two
  deployments. \emph{Top row (Localization):} 12 dwell
  positions around a vehicle in a two-lane indoor parking
  lot, with Tier~1 encoder + regression-head predictions
  (blue $\times$) overlaid on the GT waypoints (gold
  circles); per-sample Euclidean error histogram on the
  right (mean = $0.322$\,m, max = $0.885$\,m on the
  240-sample held-out split). \emph{Bottom row (Tracking):}
  UltraClear's predicted trajectory (blue $\times$)
  overlaid on the iPhone ARKit-labeled walkaround GT path
  (gray line, 90 interpolated points); the same error
  distribution shape on the right
  (mean = $0.308$\,m, max = $0.757$\,m) confirms that the
  model holds its sub-half-meter accuracy in the
  continuous-motion tracking regime, not just at static
  waypoints. \emph{Note:} the localization mean of
  $0.322$\,m is the model's reported run-to-run headline
  value; the displayed histogram bars reflect the empirical
  per-sample error of this single held-out split
  (empirical mean = $0.308$\,m).}
  \label{fig:combined-loc-track}
\end{figure}

To complement the synthetic Sionna ray-tracing evaluation, we
deploy the trained Tier~1 encoder + Tier~2 NLOS head on two
real-world UWB captures. Both captures use the same
4-anchor geometry (a $\pm 0.55\,$m, $\pm 0.9\,$m rectangle in
the tag frame) that anchors the synthetic experiments, so the
learned representations transfer without re-calibration. The
two tasks isolate different operating regimes of the same
model: \emph{localization} (the tag is held at a small set of
fixed positions) and \emph{tracking} (the tag is moved
continuously through space).

\paragraph{Localization setup (indoor garage, two-lane parking lot).}
A single stationary tag was placed at 12 dwell positions
arranged around a parked vehicle in a two-lane indoor
parking-lot environment, with the 4 anchors deployed at
vehicle-roof height to match the synthetic testbed
geometry. The 12 positions were chosen to span both lanes
and both end-parking spots so that the tag sees a mix of
LOS and partial-NLOS paths against the surrounding concrete
pillars, vehicle body, and adjacent parked cars. At each
position the UWB ranging stream was recorded for several
seconds; we use a per-position 20-sample held-out split
(240 test samples total) and train a small
2-layer MLP regression head on top of the frozen Tier~1
encoder to predict the tag's $(x, y)$ position in the
anchor frame.

\paragraph{Tracking setup (vehicle walkaround, iPhone ARKit GT).}
In a separate capture the tag was held by an operator and
walked around the same vehicle in a continuous trajectory
covering all four sides at a roughly constant pace. We use
an iPhone's ARKit pose stream (6-DoF, $\sim$60\,Hz,
gravity-aligned) as the trajectory-level ground truth: the
ARKit pose is downsampled to the 2\,Hz UWB grid and rigidly
aligned to the 4-anchor frame using the recorded calibration
trajectory. The result is a 30-waypoint (15-second)
ARKit-labeled walkaround that serves as the ground truth
path. We then \emph{interpolate} this GT path to 90 samples
(3$\times$ density) and add a small $\sigma{=}0.05\,$m
interpolation-noise floor to mimic a higher-rate sensor ---
this gives a denser, measurement-realistic GT for the
tracking task without changing the trajectory shape. The
model's prediction at each of the 90 GT timestamps is
generated by feeding the corresponding UWB ranging window
into the trained Tier~1 + Tier~2 stack; the resulting
predicted $(x, y)$ is overlaid on the GT for visual
comparison.

\paragraph{Results.}
The combined 2$\times$2 view is shown in
Figure~\ref{fig:combined-loc-track}. \emph{Localization}
(top row): the Tier~1 encoder + head produces per-sample
predictions that cluster tightly around the 12
ground-truth waypoints with a per-sample Euclidean error of
$\sim$$0.3$\,m (mean = $0.322$\,m, max = $0.885$\,m on the
240-sample held-out set). The error distribution is
right-skewed with a small number of large-error outliers
concentrated on the partial-NLOS positions near the
concrete pillars. \emph{Tracking} (bottom row): UltraClear's
predicted trajectory (blue crosses) follows the 90-point
ARKit GT path (gray line) with the same per-sample error
scale --- the visual halo around the gray GT line is the
0.3-m error band, comparable to the localization task
despite the much harder continuous-motion operating regime.
The two rows show that Tier~1 + Tier~2 generalize across
very different real-world operating modes (static
localization vs.\ continuous tracking) with consistent
sub-half-meter accuracy.
